% sections/rate_people.tex: Section 5.3: the rate from the within-borrower design, across people, ages, places and periods.
% Numbers: results/rho_within_readings.csv (within_reading.R; I6f, I6h, I6i), tab_rho_within, tab_within_by_group (appendix).
\subsection{The Rate Within Borrower, Across People}\label{sec:rate_people}

The design compares two auto loans of the same borrower, within borrower and origination-year cells, on the 13.5 million borrowers with two or more auto loans in the ten-percent archive, 65 million loans. Nothing stable about the borrower can explain a difference in her default across her own loans, so the design removes selection into contract length, the confounder that Proposition \ref{prop:termmargin} names for the term margin across borrowers. For a group of borrowers, the ratio of the within-borrower term elasticity of default within twenty-four months to the within-borrower payment elasticity, $R_{24}$, is mapped to a rate through the window formula of equation \eqref{eq:window}, with the share $\omega_t$ of the window's defaults in each month measured in the panel, the group's own marginal-utility profile, the group's exit hazard net of servicer transfers (Appendix \ref{appendix:data}; transfers are 15 percent of exits) and $T = 60$ months. The profile is the function of the group's per-period default probability $p$ that Appendix \ref{appendix:mucorrection} derives under constant relative risk aversion $\gamma$, lognormal income with dispersion $s$ per period and the persistence $\varrho$ of the distress state measured in the panel,
\begin{equation}\label{eq:lambda_p}
\lambda_k(p) \;=\; \mathbb{E}\big[(y_k - r)^{-\gamma} \,\big|\, y_k > y^*(p),\, y_1 = y^*(p)\big]\,\big(y^*(p) - r\big)^{\gamma}, \qquad \log y_k = \varrho^{\,k-1}\log y^*(p) + \sqrt{1 - \varrho^{\,2(k-1)}}\; s\, z,
\end{equation}
with $y^*(p)$ the income threshold at which a share $p$ defaults, $r$ the payment's share of income and $z$ standard normal, and with $\gamma$ set so that the constant-income limit passes through the $0.20$ that Indarte's pair measures at a mortgagor's payment share and monthly bankruptcy rate, at $s = 0.2$; across the bureau's groups $\lambda_2$ is $0.88$ to $0.89$ and $\lambda_{60}$ runs from $0.26$ in the top score decile to $0.40$ in the bottom (Table \ref{tab:rho_within}). The ratio removes the scale, and the formula is the mapping the paper applies to the experiments. The formula has a ceiling, its value at a rate of zero, $0.59$ at the pooled inputs; a ratio at or above the ceiling gives a rate of zero and is a bound.

What holding the borrower fixed does not remove is the penalty channel of Proposition \ref{prop:termmargin}, which is a property of the contract and not of who selects into it: a longer contract at the same payment is a larger balance, and default forfeits more of it. The twelve-month within-borrower ratio is negative, $-0.70$ (standard error $0.02$) pooled, because the penalty channel dominates early; the deficiency contrast of Section \ref{sec:termmargin} measures that channel at its deficiency part, a term elasticity higher by $0.36$ ($0.13$) at twelve months where the balance cannot be pursued. The twenty-four-month estimate is taken with the penalty channel unsubtracted, and it is therefore stated on the condition that the channel has decayed with the balance by then. Appendix \ref{appendix:ladder} is the check: it starts from the across-borrower term margin at twelve months, adjusts for the measured penalty and the measured selection one at a time, and arrives at $0.04$, with a set that spans the grid because the corrections are measured once on pooled samples.

\paragraph{The cells and the window.} Two choices fix the design, and each follows a rule stated in advance. The cells are borrower and origination year. The variation in contract length that identifies the far-payment response is the menu the borrower faced at each origination, which differs across lenders at a date \citep{hertzberg_screening_2018, argyle_monthly_2020}, and a borrower's two loans are almost always at different lenders; the rule is to keep that variation and remove the borrower. Adding lender-by-year cells removes the menu as well and leaves only the contract chosen within one lender's menu, which is the borrower's own selection into length, the variation the design is meant to exclude: with borrower and lender-year cells the term response at twenty-four months is $0.00$ ($0.01$), a ratio of $0.005$ that maps to no weight on far payments, a rate above the grid, which the Dobbie--Song pair rejects at the five percent level; with lender-year cells and no borrower effects the ratio is $0.62$, above the ceiling; with borrower and origination-year cells it is $0.55$ (Appendix \ref{appendix:ivvalidation}). The rule, keep the menu and remove the person, is the one that selected the population cells in Section \ref{sec:identification}, and it is the one design of the three whose ratio lies inside the formula's range. The window is twenty-four months. A fixed window removes exposure, which a lifetime outcome contains; within twelve months the penalty channel dominates and the ratio is negative ($-0.70$), so no rate can be obtained there; twenty-four months is the first window at which the far payments' weight exceeds the penalty, and the rate is stated there on the condition that the penalty channel has decayed with the balance. The precision argument of Section \ref{sec:mapping} says why the horizon of an auto loan is where the rate can be estimated at all: the far payment of a five-year contract sits near the horizon at which a payment response identifies the rate most precisely, and a thirty-year contract's far payment does not.

\paragraph{The estimate.} Table \ref{tab:rho_within} reports the results by score, income, age and origination period, and Table \ref{tab:rho_within_states} and Figure \ref{fig:rho_within} in Appendix \ref{appendix:figures} add the states. By score decile from the bottom the rates are $0.14$, $0.11$, $0.10$, $0.11$, $0.17$, $0.22$, $0.12$, $0.10$, $0.05$ and $0.01$, with 95 percent sets of $\pm 0.04$ to $\pm 0.08$ in the bottom eight deciles and wider in the top two, where defaults are rare and the sets include zero. The loan-weighted mean across deciles is $0.11$, with a 95 percent set from $0.09$ to $0.13$ that combines the ten decile sets, and this is the paper's estimate: the average of the rates the population's members reveal. The pooled ratio, $0.55$ ($0.01$) on the whole sample, is a different object. It divides the population's average term response by its average payment response, and because the deciles with the largest payment responses also have the largest term responses, the ratio of the averages exceeds every decile's own ratio, which run from $0.32$ to $0.53$, and sits near the formula's ceiling: it gives $0.02$ with a set from $0.01$ to $0.04$, the rate implied by the response of the average loan, rather than an estimate of the rate of the average borrower. Two corrections separate the estimate from the date-1 formula under a constant premium, which returns $0.09$ from the same ratios, and they run in opposite directions: assigning each default to its own decision month, as equation \eqref{eq:window} does, raises the decile mean to $0.27$, because a default in the second year was decided with the far payment nearer than the date-1 formula assumes; the declining profile of the premium lowers it to $0.11$, because the near payments have more of the weight than a constant premium assigns them. Set against the experiments, the Dobbie--Song pair's $p$-value at $0.11$ is $0.73$, and the pair accepts every rate from zero to $1.2$ at the five percent level once its five-year outcome window is aggregated the same way (Section \ref{sec:quasi}): the experiment is consistent with the estimate and does not discriminate its level.

\paragraph{Across people.} There is no gradient in score or in income. The rates in the bottom four deciles, $0.10$ to $0.14$, and in the top four, $0.01$ to $0.12$, have sets that all contain the decile mean, while default rates fall more than fifty-fold across the same deciles; by tercile the rates are $0.11$, $0.15$ and $0.06$, and by income quartile $0.24$, $0.06$, $0.08$ and $0.26$, with sets that overlap or meet at $0.15$. The measured ratios differ little across deciles, $0.53$ at the bottom and $0.49$ at the top; the inputs differ more, and in the direction that lowers the top: a common ratio of $0.44$ mapped with each decile's own profile and exit hazard is $0.21$ in the bottom decile and $0.05$ in the top, because the far payments' weight falls as default becomes rare ($\lambda_{60}$ from $0.40$ to $0.26$). The table reports each decile at its own inputs, and the decile sets are what the absence of a gradient rests on. The level moves with the premium's shape: with a constant premium of $0.2$ at every later date in place of the profile the decile mean is $0.24$, and with $0.5$ it is $0.29$ (\texttt{results/rho\_within\_sensitivity.csv}). The fifth and sixth deciles, at $0.17$ and $0.22$, sit above the others, and the theory of Section \ref{sec:results} predicts no such rise under a common rate. The rise is in the denominator: the within-borrower term response is nearly flat from the third to the sixth decile, $0.27$ to $0.29$ (Table \ref{tab:within_by_group}, Appendix \ref{appendix:ladder}), while the payment response rises from $0.56$ to $0.84$, so the deciles that respond most to the payment, the exposure peak of Section \ref{sec:heterogeneity}, have the lowest ratios. Two mechanisms could produce it and the data speak to both. The penalty channel lowers the term response most where balances are charged off most, in the bottom deciles, and would therefore depress the bottom deciles' ratios below the middle's; they are not depressed, so the penalty channel does not produce the rise. Payoff expectations could: a borrower who expects to refinance or trade in sooner than the measured hazard weighs the far payment less and reveals a higher rate (Proposition \ref{prop:beliefs}), and a gradient in those expectations across the middle of the score distribution is the one belief the survey of Appendix \ref{appendix:beliefs} cannot bound. The paper reports the middle deciles as measured. The experiments' subgroup pairs, which assign the timing of payments within one distressed population, find the same rate across credit score, leverage and employment, with equality not rejected at $p = 0.84$, $0.64$ and $0.21$ (Table \ref{tab:rho_subgroups}).

\paragraph{Age.} By age the rates are $0.32$ for borrowers aged 18 to 30, with a 95 percent set from $0.21$ to $0.48$, and zero at every older age, with sets that end at $0.10$ for ages 31 to 40, $0.03$ for 41 to 50, zero for 51 to 60 and $0.07$ above 60. The youngest group's ratio is $0.26$ against $0.65$ above 60; the payment elasticity rises with age, from $0.54$ to $0.92$, and the term elasticity rises faster, from $0.14$ to $0.60$. Age is the one dimension along which the within-borrower design finds a gradient.

\paragraph{Places and periods.} Across the twelve largest states the ratio ranges from $-0.21$ in California to $1.20$ in Texas, while the formula's range over every rate from zero to infinity is zero to $0.59$ at the pooled inputs; six states sit at or above the ceiling, five inside the range (Michigan $0.02$, Florida $0.10$, Pennsylvania $0.14$, New York $0.18$, New Jersey $0.21$), and California's term response is negative at twenty-four months, as every state's is at twelve. A ratio outside the formula's range in either direction is the term margin measuring something other than the weight on far payments, and Proposition \ref{prop:termmargin} names the penalty channel, whose size differs across states with their collection law (Section \ref{sec:termmargin}); the paper therefore does not assign a rate to a place from the term margin, and the results across places in Section \ref{sec:heterogeneity} are on the payment margin. By origination period, from one regression on the whole within-borrower sample with the borrower effect identified across years and the two elasticities allowed to differ by the loan's period, the ratio at twenty-four months falls from $0.75$ for loans opened in 2005 to 2009 and $0.59$ for 2010 to 2013, both at or above the ceiling, to $0.54$ for 2014 to 2017, a rate of $0.04$ with a set from $0.02$ to $0.06$, and $0.38$ for 2018 to 2019, a rate of $0.18$ with a set from $0.14$ to $0.22$; for loans opened in 2020 to 2022 the ratio at twenty-four months is negative, $-0.05$ ($0.03$), because the term response is $-0.02$, and there is no rate. The decline is in the term elasticity, from $0.39$ to $-0.02$, while the payment elasticity moves between $0.44$ and $0.66$ without trend, and the paper does not interpret it as a trend in the rate: the later periods' windows end closer to the archive's last month, the mix of contract lengths shifted toward 72 and 84 months over the period, and the pandemic cohorts' negative term response has the sign of the penalty channel.

\begin{table}[H]
\centering
\caption{The rate across people from the within-borrower design}
\label{tab:rho_within}
\small
\input{tables/tab_rho_within_main}
\vspace{0.5em}
\begin{minipage}{0.92\textwidth}
\noindent \small \textit{Notes:} $R_{24}$ is the within-borrower term elasticity of default within twenty-four months over the within-borrower payment elasticity, on borrowers with two or more auto loans in the ten-percent archive (13.5 million borrowers, 65 million loans; Table \ref{tab:within_by_group}), with the delta-method standard error. $\lambda_2$ and $\lambda_{60}$ are the group's marginal-utility profile of equation \eqref{eq:lambda_p} at the group's default rate, at the second and the sixtieth payment; $h$ is the group's exit hazard net of servicer transfers (Appendix \ref{appendix:data}), for score deciles the exposure-weighted mean over contract-length bins; $T = 60$ months. $\rho$ solves equation \eqref{eq:window} at $K = 24$ with the panel's default timing; a ratio at or below zero has no rate (--); a ratio at or above the formula's value at $\rho = 0$ gives zero. The set carries the ratio's standard error through the formula to the rate. The twelve largest states are in Table \ref{tab:rho_within_states}.
\end{minipage}
% replication: Code/identification/I6f_within_person_by_group.R, Code/identification/I6i_within_person_by_period.R, Code/local/within_reading.R
\end{table}
